% ============================================================
%  Příprava na maturitu – Mechatronika
%  Kompiluj: pdflatex maturita_mechatronika.tex
%  Kódování: UTF-8
% ============================================================

\documentclass[a4paper,10pt]{article}

% --- Balíčky ---
\usepackage[T1]{fontenc}
\usepackage[utf8]{inputenc}
\usepackage[czech]{babel}
\usepackage[left=2.4cm, right=1.8cm, top=1.8cm, bottom=1.8cm]{geometry}
\usepackage{lmodern}
\usepackage{microtype}
\usepackage{parskip}
\usepackage{titlesec}
\usepackage{enumitem}
\usepackage{xcolor}
\usepackage{hyperref}
\usepackage{tabularx}
\usepackage{booktabs}
\usepackage{array}
\usepackage{tikz}
\usepackage{eso-pic}
\usepackage{mdframed}
\usepackage{amsmath}
\usepackage{amssymb}
\usepackage{pgfplots}
\pgfplotsset{compat=1.18}
\usetikzlibrary{arrows.meta, positioning, shapes.geometric, calc}

% --- Barvy ---
\definecolor{accent}{HTML}{03254E}
\definecolor{stripeblue}{HTML}{568EA3}
\definecolor{medgray}{HTML}{888888}
\definecolor{lightblue}{HTML}{EAF2F8}
\definecolor{lightyellow}{HTML}{FDFBE4}
\definecolor{lightgreen}{HTML}{EAF7EA}

% --- Modrý pruh vlevo ---
\AddToShipoutPictureBG{%
  \begin{tikzpicture}[remember picture, overlay]
    \fill[stripeblue]
      (current page.north west)
      rectangle
      ([xshift=10mm] current page.south west);
    \fill[accent!60]
      ([xshift=10mm] current page.north west)
      rectangle
      ([xshift=12mm] current page.south west);
  \end{tikzpicture}%
}

\hypersetup{colorlinks=true, urlcolor=accent, linkcolor=accent}

\titleformat{\section}
  {\large\bfseries\color{accent}}
  {\thesection.}
  {0.5em}
  {}
  [\color{accent}\titlerule]

\titlespacing{\section}{0pt}{16pt}{6pt}

\newmdenv[
  backgroundcolor=lightblue,
  linecolor=stripeblue,
  linewidth=1pt,
  roundcorner=4pt,
  innerleftmargin=8pt,
  innerrightmargin=8pt,
  innertopmargin=6pt,
  innerbottommargin=6pt,
  skipabove=4pt,
  skipbelow=4pt
]{pojmybox}

\newmdenv[
  backgroundcolor=lightyellow,
  linecolor=accent!40,
  linewidth=1pt,
  roundcorner=4pt,
  innerleftmargin=8pt,
  innerrightmargin=8pt,
  innertopmargin=6pt,
  innerbottommargin=6pt,
  skipabove=4pt,
  skipbelow=8pt
]{vysvetlenibox}

\newmdenv[
  backgroundcolor=lightgreen,
  linecolor=accent!30,
  linewidth=1pt,
  roundcorner=4pt,
  innerleftmargin=8pt,
  innerrightmargin=8pt,
  innertopmargin=6pt,
  innerbottommargin=6pt,
  skipabove=4pt,
  skipbelow=8pt
]{vzorcebox}

\pagestyle{plain}

% ============================================================
\begin{document}

% --- Záhlaví ---
\begin{minipage}[t]{0.75\textwidth}
    {\Huge\bfseries\color{accent} Příprava na maturitu}\\[4pt]
    {\large\color{stripeblue} Mechatronika}\\[2pt]
    \textcolor{medgray}{\small Suchánek Lukáš \quad SPŠ Prosek \quad 2026}
\end{minipage}
\hfill
\begin{minipage}[t]{0.22\textwidth}
    \raggedleft
    \begin{tikzpicture}
        \draw[color=stripeblue, rounded corners=4pt, line width=1pt]
            (0,0) rectangle (3,1.2)
            node[midway, align=center, color=accent, font=\small\bfseries]
            {Otázek: 22};
    \end{tikzpicture}
\end{minipage}

\vspace{6pt}
\noindent\color{accent}\rule{\linewidth}{1.5pt}
\vspace{2pt}

% ============================================================
\section{Kombinační logické řízení, minimalizace}

\begin{pojmybox}
\textbf{Klíčové pojmy:}
kombinační obvod, DNF, KNF, Karnaughova mapa, minterm, maxterm,
Quine-McCluskey, multiplexer, demultiplexer, dekodér, enkodér, sčítačka
\end{pojmybox}

\begin{vysvetlenibox}
\textbf{Kombinační obvod}\\
Výstup závisí \textit{pouze} na okamžitých vstupech -- obvod nemá paměť.
Logická funkce se popisuje pravdivostní tabulkou nebo algebraickým výrazem.

\textbf{Příklad: Hlasovací systém}\\
Mějme 3 členy komise (A, B, C). Návrh projde, když hlasují \textit{alespoň 2} pro.
\begin{center}
\begin{tabular}{ccc|c}
    \toprule
    A & B & C & Výstup Y \\
    \midrule
    0 & 0 & 0 & 0 \\
    0 & 0 & 1 & 0 \\
    0 & 1 & 0 & 0 \\
    0 & 1 & 1 & 1 \\
    1 & 0 & 0 & 0 \\
    1 & 0 & 1 & 1 \\
    1 & 1 & 0 & 1 \\
    1 & 1 & 1 & 1 \\
    \bottomrule
\end{tabular}
\end{center}

\textbf{DNF a KNF}\\
Postup vytvoření normálních forem z pravdivostní tabulky:
\begin{itemize}[noitemsep, leftmargin=1.4em, topsep=4pt]
    \item \textbf{DNF} (Disjunktivní normální forma) -- součet součinů
          (mintermů); pro každý řádek, kde výstup = 1.
          \textit{Příklad:} $Y = \overline{A}BC + A\overline{B}C + AB\overline{C} + ABC$
    \item \textbf{KNF} (Konjunktivní normální forma) -- součin součtů
          (maxtermů); pro každý řádek, kde výstup = 0.
          \textit{Příklad:} $Y = (A+B+C)(A+B+\overline{C})(A+\overline{B}+C)(\overline{A}+B+C)$
\end{itemize}

\textbf{Postup minimalizace pomocí Karnaughovy mapy:}
\begin{enumerate}[noitemsep, leftmargin=1.4em, topsep=4pt]
    \item Sestrojíme K-mapu z pravdivostní tabulky
    \item Seskupíme jedničky do skupin o velikosti $2^n$ (1, 2, 4, 8\ldots)
    \item Skupiny mohou přesahovat přes okraje (mapa je "srolovaná")
    \item Z každé skupiny vyvodíme zjednodušený výraz
    \item Výsledná funkce je součtem všech skupin
\end{enumerate}

\textbf{Karnaughova mapa -- postup minimalizace}\\
Grafická metoda pro minimalizaci až 4 proměnných.
Sousední buňky se liší v \textit{jedné} proměnné (Grayův kód).
Seskupujeme jedničky do skupin o velikosti $2^n$ (1, 2, 4, 8\ldots).
Čím větší skupina, tím jednodušší výsledný výraz.

\textbf{Pravidla pro Karnaughovu mapu:}
\begin{itemize}[noitemsep, leftmargin=1.4em, topsep=4pt]
    \item \textbf{Počet buněk:} $2^n$ kde $n$ je počet proměnných (2→4, 3→8, 4→16)
    \item \textbf{Číslování os:} Grayův kód (00, 01, 11, 10) -- sousedi se liší o 1 bit
    \item \textbf{Seskupování:} skupiny musí být obdélníkové, velikost $2^n$
    \item \textbf{Přesahy:} mapa je "nekonečná" -- levý okraj sousedí s pravým
    \item \textbf{Cíl:} co nejméně skupin, každá co největší
\end{itemize}

\begin{center}
\begin{tikzpicture}[scale=0.75]
  % Karnaughova mapa 3 proměnné (A, BC)
  \draw[accent, thick] (0,0) grid (4,2);
  % Záhlaví
  \node at (-0.6, 1.5) {\small $A=0$};
  \node at (-0.6, 0.5) {\small $A=1$};
  \node at (0.5, 2.3)  {\small $BC{=}00$};
  \node at (1.5, 2.3)  {\small $BC{=}01$};
  \node at (2.5, 2.3)  {\small $BC{=}11$};
  \node at (3.5, 2.3)  {\small $BC{=}10$};
  % Hodnoty (příklad: f = A·B + C)
  \node at (0.5,1.5) {0};
  \node at (1.5,1.5) {1};
  \node at (2.5,1.5) {1};
  \node at (3.5,1.5) {0};
  \node at (0.5,0.5) {0};
  \node at (1.5,0.5) {1};
  \node at (2.5,0.5) {1};
  \node at (3.5,0.5) {1};
  % Skupiny
  \draw[stripeblue, very thick, rounded corners=4pt]
    (1.08,0.08) rectangle (2.92,1.92);
  \draw[red!70!black, very thick, rounded corners=4pt]
    (2.08,-0.05) rectangle (4.05,0.95);
  % Popisky skupin
  \node[color=stripeblue, font=\scriptsize] at (2,2.7) {skupina: $B$};
  \node[color=red!70!black, font=\scriptsize] at (3.5,-0.4) {skupina: $A\cdot C$};
\end{tikzpicture}
\end{center}

\begin{vzorcebox}
Výsledek minimalizace: $f = B + A \cdot \overline{C}$
\quad (místo původního rozsáhlého DNF výrazu)
\end{vzorcebox}

\textbf{Quine-McCluskeyho metoda}\\
Pro více než 4 proměnné se používá algoritmická metoda:
\begin{enumerate}[noitemsep, leftmargin=1.4em, topsep=4pt]
    \item Seřadíme mintermy podle počtu jedniček v binárním zápisu
    \item Hledáme dvojice, které se liší v jednom bitu a slučujeme je
    \item Opakujeme, dokud nelze dále slučovat
    \item Z výsledných prvočlenů vybereme minimální pokrytí
\end{enumerate}
Tato metoda je vhodná pro počítačovou implementaci.

\textbf{Typické kombinační obvody}

\begin{center}
\begin{tabularx}{0.95\linewidth}{@{} l X @{}}
    \toprule
    \textbf{Obvod}    & \textbf{Funkce} \\
    \midrule
    Multiplexer (MUX) & Vybírá 1 ze $2^n$ vstupů a přepošle na výstup dle adresních vodičů \\
    Demultiplexer     & Směruje 1 vstup na 1 z $2^n$ výstupů \\
    Dekodér           & Převádí $n$-bitový kód na $2^n$ výstupů (aktivní vždy jen 1) \\
    Enkodér           & Opak dekodéru -- $2^n$ vstupů na $n$-bitový kód \\
    Polosčítačka      & Sčítá 2 bity, výstupy: součet $S$ a přenos $C$ \\
    Úplná sčítačka    & Sčítá 2 bity + vstupní přenos; základ ALU \\
    Komparátor        & Porovnává dvě $n$-bitová čísla ($A>B$, $A=B$, $A<B$) \\
    \bottomrule
\end{tabularx}
\end{center}

\textbf{Multiplexer v praxi}\\
Multiplexer funguje jako "číselný přepínač". Příklad: 4:1 MUX má 4 datové vstupy
($D_0$--$D_3$), 2 adresní vstupy ($A_0$, $A_1$) a 1 výstup ($Y$).
\begin{itemize}[noitemsep, leftmargin=1.4em, topsep=4pt]
    \item Pokud $A_1A_0 = 00$, na výstup jde $D_0$
    \item Pokud $A_1A_0 = 01$, na výstup jde $D_1$
    \item Pokud $A_1A_0 = 10$, na výstup jde $D_2$
    \item Pokud $A_1A_0 = 11$, na výstup jde $D_3$
\end{itemize}
\textit{Použití:} výběr signálu z více zdrojů, paralelně-sériový převod.

\textbf{Dekodér a enkodér}\\
\begin{itemize}[noitemsep, leftmargin=1.4em, topsep=4pt]
    \item \textbf{Dekodér 3:8} -- 3 bity na vstupu aktivují 1 z 8 výstupů.
          Použití: adresa paměti, výběr periférie, řízení 7-segmentového displeje.
    \item \textbf{Enkodér 8:3} -- 8 vstupů, aktivní vždy jen jeden; výstup je 3bitový kód.
          Prioritní enkodér řeší situaci, kdy je aktivních více vstupů (vybere nejvyšší).
\end{itemize}

\textbf{Sčítačky a ALU}\\
Polosčítačka sčítá 2 bity, ale neumí přenos ze sčítání nižšího řádu.
\begin{itemize}[noitemsep, leftmargin=1.4em, topsep=4pt]
    \item \textbf{Polosčítačka} -- 2 vstupy ($A$, $B$), 2 výstupy ($S$, $C_{out}$)
    \item \textbf{Úplná sčítačka} -- 3 vstupy ($A$, $B$, $C_{in}$), 2 výstupy ($S$, $C_{out}$)
    \item \textbf{4bitová sčítačka} -- 4 úplné sčítačky za sebou; $C_{out}$ z předchozí jde do $C_{in}$ další
\end{itemize}
ALU (Arithmetic Logic Unit) v procesoru obsahuje sčítačky, logické funkce
a posuvné registry -- vše jsou kombinační obvody.

\begin{vzorcebox}
Polosčítačka: $S = A \oplus B$, \quad $C = A \cdot B$\\
Úplná sčítačka: $S = A \oplus B \oplus C_{in}$, \quad
$C_{out} = A \cdot B + C_{in} \cdot (A \oplus B)$\\
4bitová sčítačka: $C_{out}$ z bitu 0 jde do $C_{in}$ bitu 1 atd.
\end{vzorcebox}

\textbf{Komparátor}\\
Porovnává dvě čísla $A$ a $B$. Výstupy: $A>B$, $A=B$, $A<B$.
\begin{itemize}[noitemsep, leftmargin=1.4em, topsep=4pt]
    \item Rovnost: $A=B$ když všechny bity souhlasí ($XNOR$ každého páru bitů)
    \item Větší: $A>B$ když nejvyšší neshodný bit je u $A$ roven 1
\end{itemize}
\textit{Použití:} řízení procesů, bezpečnostní limity, třídění dat.
\end{vysvetlenibox}

% ============================================================

% ============================================================
\end{document}
